\documentclass[11pt,a4paper]{article}

\usepackage[utf8]{inputenc}
\usepackage[T1]{fontenc}
\usepackage{amsmath,amssymb,amsfonts}
\usepackage{geometry}
\usepackage{cite}
\usepackage{booktabs}

\geometry{margin=1in}

\title{Project Abstract for AWS: G-quAI \\ \large Hardware-Native Quantum Convolutional Neural Networks on QuEra Aquila for Medical Image Classification}
\author{Team G-quAI, University of Coimbra}
\date{\today}

\begin{document}

\maketitle

\begin{abstract}
\textbf{G-quAI} is a research initiative based at the University of Coimbra dedicated to the early detection of colorectal cancer from high-resolution gastrointestinal medical imaging using Quantum Artificial Intelligence. Classical deep learning for such high-variance medical imaging requires massive computational clusters. We aim to leverage hardware-native quantum phenomena to reduce this computational overhead while increasing diagnostic sensitivity. While our ultimate goal targets colorectal cancer, this specific experimental proposal serves as a crucial hardware proof-of-concept utilizing the \textbf{BreastMNIST} dataset. BreastMNIST contains highly noisy, low-resolution ultrasound images with a manageable low sample count, making it an ideal, iterative testbed for quantifying the diagnostic advantage of quantum physics prior to scaling to high-resolution endoscopic datasets. For this phase, we are requesting resources to run our \textit{Digital} encoding experiments on the QuEra Aquila system.
\end{abstract}

\section*{1. Motivation: Why the Quantum Domain?}
Our work is grounded in the architectural advantages of Digital-Analog Quantum Kernels (DAQKs) as proposed by Simen et al. \cite{daqcnn2024}. The transition to the quantum domain for medical image classification is motivated by three primary factors:
\begin{enumerate}
    \item \textbf{High-Dimensional Feature Mapping:} Quantum kernels naturally map classical data into a high-dimensional Hilbert space where classes that are non-linearly separable in the original space may become linearly separable.
    \item \textbf{Complexity via Native Interactions:} Unlike classical CNNs where filters are trained via backpropagation, DAQCNNs leverage the fixed, native "complexity" of Rydberg atom interactions to extract intricate spatial correlations and multipartite entanglement dynamics that classical filters may fail to capture.
    \item \textbf{Massive Parameter Efficiency:} DAQCNNs have demonstrated the ability to outperform classical state-of-the-art models while using significantly fewer trainable parameters (often by orders of magnitude), directly addressing the energy and computational bottlenecks of modern medical AI.
\end{enumerate}

\section*{2. Methodology: ZZ-AQCNN \& Encoding Schemes}
Our approach, the \textbf{ZZ-AQCNN}, enhances the DAQCNN framework by measuring both 1-body ($\langle \sigma_Z^{(i)} \rangle$) and 2-body ($\langle \sigma_Z^{(i)} \sigma_Z^{(j)} \rangle$) correlators. For a 9-qubit kernel, this expands our feature space to 45 densely correlated features. 

We utilize two primary encoding schemes:
\begin{enumerate}
    \item \textbf{Digital Encoding (Requested):} Data is encoded in the initial state via $RY(x_i)$ and $H$ gates at $t=0$:
    \begin{equation}
        |\Psi_{out}\rangle = e^{-i H_{fixed} T} \left( \bigotimes_{i=1}^N H \cdot RY(x_i) \right) |0\rangle^{\otimes N}
    \end{equation}
    
    \item \textbf{Analog Encoding:} Data is injected directly into the Hamiltonian dynamics via local detuning:
    \begin{equation}
        |\Psi_{out}\rangle = \mathcal{T} \exp \left( -i \int_0^T H(\mathbf{x}, t) dt \right) |0\rangle^{\otimes N}
    \end{equation}
\end{enumerate}
\textbf{Scope:} We focus on \textbf{Digital Encoding} for this request to establish a robust hardware baseline.

\section*{3. Preliminary Results (Trotterized Simulation)}
Extensive simulations on BreastMNIST using Trotterized Hamiltonian evolution ($T=2.5, dt=0.05$) have yielded the following performance metrics across 10 independent validation seeds:

\begin{table}[htbp]
    \centering
    \begin{tabular}{l c c c c}
        \toprule
        Config / Metric & Classical & Original DAQCNN & Digital + ZZ & Analog + ZZ \\
                        & Baseline & (No ZZ) \cite{daqcnn2024} & (Enhanced) & (ZZ-AQCNN) \\
        \midrule
        \textbf{Test Acc (\%)} & $0.0 \pm 0.0$ & $86.9 \pm 1.2$ & $87.2 \pm 0.9$ & $86.6 \pm 1.0$ \\
        \textbf{ROC-AUC}       & $0.00 \pm 0.0$ & $0.88 \pm 0.01$ & $0.91 \pm 0.01$ & $0.93 \pm 0.01$ \\
        \bottomrule
    \end{tabular}
    \caption{Simulated results on BreastMNIST. Note the superior separability (AUC) of the ZZ-enhanced models.}
\end{table}

Feature importance analysis confirms that \textbf{79.5\% of the predictive power} stems from the $ZZ$ correlators, validating our architectural expansion.

\section*{4. Experimental Plan \& Resource Estimation}
We utilize \textbf{Spatial Multiplexing} on the QuEra Aquila (256 qubits) to process 21 patches ($3\times3$) per task. Quantum results will be \textbf{cached} for offline training of the classical head.

\textbf{Workload:} 780 images $\times$ 81 patches / 21 parallel $\approx$ \textbf{3,009 tasks}.

\begin{table}[h]
    \centering
    \begin{tabular}{lrrr}
        \toprule
        Precision ($\epsilon$) & Shots per Task & Total Shots & \textbf{Estimated Cost (\$0.01/shot)} \\
        \midrule
        0.2 (20\% Error) & 25 & 75,225 & \textbf{\$752.25} \\
        0.1 (10\% Error) & 100 & 300,900 & \textbf{\$3,009.00} \\
        0.05 (5\% Error) & 400 & 1,203,600 & \textbf{\$12,036.00} \\
        \bottomrule
    \end{tabular}
\end{table}

\begin{thebibliography}{9}
\bibitem{daqcnn2024}
Simen, A., Flores-Garrigos, C., Hegade, N. N., et al. (2024). \textit{Digital-analog quantum convolutional neural networks for image classification}. Phys. Rev. Research, 6(4), L042060.
\end{thebibliography}

\end{document}
